
\documentclass[10pt]{article}
\usepackage[T1]{fontenc}
\usepackage[utf8]{inputenc}
\usepackage{lmodern}
\usepackage{amsmath,amssymb,amsthm,mathrsfs,mathtools}
\usepackage{geometry}
\usepackage{enumitem}
\usepackage{tabularx}
\geometry{margin=0.82in}
\setlength{\parindent}{1.5em}
\setlength{\parskip}{0.16em}
\emergencystretch=2em
\setlist{leftmargin=2.0em,itemsep=0.1em,topsep=0.2em}
\newcommand{\A}{\mathbb A}
\newcommand{\Af}{\mathbb A_f}
\newcommand{\C}{\mathbb C}
\newcommand{\G}{\mathbb G}
\newcommand{\Q}{\mathbb Q}
\newcommand{\R}{\mathbb R}
\newcommand{\Z}{\mathbb Z}
\newcommand{\bZ}{\widehat{\mathbb Z}}
\newcommand{\Ga}{\mathbb G_a}
\newcommand{\Gm}{\mathbb G_m}
\newcommand{\GL}{\operatorname{GL}}
\newcommand{\SL}{\operatorname{SL}}
\newcommand{\Gal}{\operatorname{Gal}}
\newcommand{\Hom}{\operatorname{Hom}}
\newcommand{\Isom}{\operatorname{Isom}}
\newcommand{\Lie}{\operatorname{Lie}}
\newcommand{\Sym}{\operatorname{Sym}}
\newcommand{\Spec}{\operatorname{Spec}}
\newcommand{\End}{\operatorname{End}}
\newcommand{\Tr}{\operatorname{Tr}}
\newcommand{\cO}{\mathcal O}
\newcommand{\cH}{\mathcal H}
\newcommand{\cU}{\mathcal U}
\newcommand{\cL}{\mathcal L}
\newcommand{\dd}{\,d}
\newcommand{\et}{\mathrm{\acute et}}
\newcommand{\an}{\mathrm{an}}
\newcommand{\coker}{\operatorname{coker}}
\newcommand{\im}{\operatorname{im}}
\newcommand{\sh}{\operatorname{sh}}
\newcommand{\QED}{\hfill$\square$}
\newcommand{\F}{\mathbb F}
\newcommand{\Imop}{\operatorname{Im}}
\newcommand{\charop}{\operatorname{char}}
\newtheorem*{proposition}{Proposition}
\newtheorem*{theorem}{Théorème}
\newtheorem*{definition}{Définition}
\newtheorem*{lemma}{Lemme}
\begin{document}

\begin{center}
{\Large\bfseries FORMES MODULAIRES ET REPRÉSENTATIONS $\ell$-ADIQUES}\par
\vspace{0.35em}
{\large Séminaire Bourbaki, 21e année, 1968/69, no. 355}\par
\vspace{0.7em}
{\itshape Pierre Deligne}
\end{center}

\section{Introduction}

Soit
\[
D(q)=q\prod_{n=1}^{\infty}(1-q^n)^{24}
     =\sum_{n=1}^{\infty}\tau(n)q^n
\qquad (|q|<1)
\]
et
\[
\Delta(z)=D(e^{2\pi iz})\qquad (\operatorname{Im} z>0).
\]
On sait que la fonction $\Delta$ est, à un facteur constant près, l'unique forme parabolique
de poids $12$ pour le groupe $\SL_2(\Z)$. Posons, pour $p$ premier,
\[
H_p(X)=1-\tau(p)X+p^{11}X^2 .
\]
D'après la théorie de Hecke, la série de Dirichlet
\[
L_\tau(s)=\sum \tau(n)n^{-s}=\prod_p\frac{1}{H_p(p^{-s})}
\]
se prolonge en une fonction entière de $s$, et la fonction
\[
(2\pi)^{-s}\Gamma(s)L_\tau(s)
\]
est invariante par $s\mapsto 12-s$. La conjecture de Ramanujan affirme que les racines
du polynôme $H_p$ sont de valeur absolue $p^{-11/2}$, c'est-à-dire que
$|\tau(p)|<2p^{11/2}$.

Ces propriétés, démontrées ou conjecturales, ressemblent aux propriétés conjecturales
des fonctions zêta des variétés algébriques sur $\Q$. Suggéré par cela, en première
approximation, on est tenté d'interpréter $L_\tau$ comme la fonction zêta d'une telle
variété.

Pour chaque nombre premier $\ell$, soit $K_\ell$ l'extension maximale de $\Q$ non
ramifiée hors de $\ell$ et, pour $p\ne\ell$, soit $F_p$ l'inverse, dans $\Gal(K_\ell/\Q)$,
de l'élément de Frobenius $\varphi_p$ relatif à $p$. Ce dernier objet est bien défini à
conjugaison près.

En traduisant ce qui précède en termes de cohomologie $\ell$-adique, Serre a conjecturé
l'existence, pour chaque $\ell$, d'une représentation de $\Gal(K_\ell/\Q)$ dans un
$\Q_\ell$-espace vectoriel $W_\ell$ de rang $2$ telle que, pour chaque $p\ne\ell$,
\[
H_p(X)=\det(1-F_pX;W_\ell).
\]
De plus, la représentation $W_\ell$ devrait relever du domaine d'application des
conjectures de Weil, et la conjecture de Ramanujan devrait en être un cas particulier.

Ce programme a été mené à bien par Kuga--Shimura \([4]\) dans le cas analogue des
formes modulaires relatives à certains sous-groupes de $\SL_2(\R)$ à quotient compact.
Réduite au cas présent, l'idée fondamentale de Sato--Kuga--Shimura est la suivante :
si $E$ est la courbe elliptique universelle sur le schéma de modules $S$ des courbes
elliptiques (en oubliant qu'il n'existe pas), et si $E^k$ est le produit fibré de $E$ avec
lui-même sur $S$, répété $k$ fois, alors $L_\tau(s)$ est essentiellement la fonction zêta
de $E^k$ pour
$k=10=12-2$.

Nous montrons ci-dessous comment résoudre les difficultés créées par les points à
l'infini, et comment construire les représentations $W_\ell$ ayant les propriétés indiquées
ci-dessus. Pour des détails historiques plus complets et pour les applications, voir Serre
\([6]\).

\paragraph{Notations.}
\begin{itemize}
\item On note $\A$ l'anneau des adèles de $\Q$, $\Af$ l'anneau des adèles finies, produit
restreint, pris sur tous les nombres premiers, des corps $\Q_p$, et, pour $S$ un ensemble
fini de nombres premiers, on pose
\[
\Af^S=\prod_{p\in S}\Q_p\times\prod_{p\notin S}\Z_p\subset \Af .
\]
Pour $S=\varnothing$, on écrit $\bZ=\Af^\varnothing$.
\item Si $X$ est un espace topologique, ou le site étale d'un schéma, et si $G$ est un
ensemble, $G$ désigne aussi le faisceau constant sur $X$ défini par $G$.
\item On note $\Ga$ et $\Gm$ les groupes additif et multiplicatif.
\item Une courbe elliptique est une variété abélienne de dimension un; en particulier, elle
est munie d'une origine.
\item Si $L$ est un faisceau inversible et si $n\in\Z$, $L^n$ désigne la $n$-ième puissance
tensorielle $L^{\otimes n}$.
\item On note $\overline{\Q}$ la clôture algébrique de $\Q$ dans $\C$.
\item Le symbole $\square$ marque la fin d'une démonstration, ou son absence.
\end{itemize}

\section{L'isomorphisme de Shimura}

\paragraph{(2.1)}
Une courbe elliptique sur un espace analytique complexe $S$ est un morphisme propre et
plat d'espaces analytiques $f:E\to S$, muni d'une section $e$, dont les fibres sont des
courbes elliptiques. Une courbe elliptique sur $S$ admet une et une seule loi de groupe sur
$S$,
\[
\mu:E\times_S E\to E,
\]
dont la section unité est $e$. À une courbe elliptique sont associés les objets suivants.

\begin{enumerate}[label=(\alph*)]
\item Le faisceau inversible
\[
\omega_E=e^*\Omega^1_{E/S}.
\]
L'algèbre de Lie relative $\Lie_S(E)$ est le faisceau inversible $\omega^{-1}$ dual de
$\omega$. On a
\[
f_*\Omega^1_{E/S}\xrightarrow{\sim}\omega .
\]

\item Le système local de $\Z$-modules libres de rang $2$, $R^1f_*\Z$. On pose
\[
T_\Z(E)=(R^1f_*\Z)^\vee,\qquad
T_\Q(E)=T_\Z(E)\otimes\Q,
\]
le système local d'homologie de $E$ sur $S$.
\end{enumerate}

L'application exponentielle définit une suite exacte de faisceaux de sections
\[
0\longrightarrow T_\Z(E)\xrightarrow{\alpha}\omega^{-1}\longrightarrow E\longrightarrow 0,
\]
de sorte que la courbe elliptique $E$ peut être reconstruite à partir de l'application
$\alpha$.

Le système local $\bigwedge^2 R^1f_*\Z\simeq R^2f_*\Z$ est canoniquement isomorphe
à $\Z$. Un isomorphisme entre $\Z^2$ et $R^1f_*\Z$ est dit admissible s'il induit $1$ sur
les secondes puissances extérieures. Notons $\Hom^+(\R^2,\C)$ l'ensemble des
isomorphismes d'espaces vectoriels réels entre $\R^2$ et $\C$ qui respectent les orientations
naturelles de $\R^2$ et de $\C$, définies par $e_1\wedge e_2>0$ et $1\wedge i>0$.
Un tel homomorphisme est déterminé par sa restriction à $\Z^2$, et l'on pose
\[
\Hom^+(\Z^2,\C)=\Hom^+(\R^2,\C).
\]
Cet espace est muni de la structure complexe obtenue par son inclusion dans l'espace
vectoriel complexe $\Hom(\Z^2,\C)$. Sur cet espace, on dispose d'une suite exacte
« universelle »
\[
0\longrightarrow \Z^2\xrightarrow{\alpha}\Ga\longrightarrow E_0\longrightarrow 0.
\]

\begin{proposition}[2.2]
\leavevmode
\begin{enumerate}[label=(\roman*)]
\item Le foncteur qui associe à tout espace analytique $S$ l'ensemble des classes
d'isomorphisme de courbes elliptiques $E$ sur $S$, munies d'isomorphismes
\[
\omega_E\simeq \Ga,\qquad R^1f_*\Z\simeq \Z^2
\]
(le dernier étant admissible), est représenté par l'espace analytique
$\Hom^+(\R^2,\C)$, muni d'une courbe elliptique universelle $E_0$.

\item Le foncteur qui associe à tout espace analytique $S$ l'ensemble des classes
d'isomorphisme de courbes elliptiques sur $S$, munies d'un isomorphisme admissible
$R^1f_*\Z\simeq \Z^2$, est représenté par l'espace analytique
\[
X=\C^\times\backslash\Hom^+(\R^2,\C),
\]
le demi-plan de Poincaré.

\item L'espace $\Hom^+(\R^2,\C)$ est un espace homogène principal du groupe $\Gm$ sur
$X$.
\end{enumerate}
\end{proposition}

On peut de nouveau regarder $X$ comme l'ensemble des structures complexes sur $\R^2$.
Cet espace est, d'après (ii), muni d'une courbe elliptique universelle $E_X$, dont le système
local de cohomologie réelle est canoniquement isomorphe à $\R^2$. Soit $\omega$ le
faisceau inversible associé à $E_X$.

Le faisceau analytique cohérent
\[
R^1f_*\R\otimes_\R\cO_X
\]
est le faisceau de cohomologie de de Rham relative de $E_X$ sur $X$, et s'insère à ce titre
dans une suite exacte, la filtration de Hodge,
\[
0\longrightarrow \omega\longrightarrow R^1f_*\R\otimes_\R\cO_X
\xrightarrow{q}\omega^{-1}\longrightarrow 0,
\]
puisque, par dualité de Serre, $\omega^{-1}\simeq R^1f_*\cO_{E_X}$.

La description fonctorielle de 2.2(ii) donne évidemment une action à droite du groupe
$\SL_2(\Z)$ sur $(X,E_X)$ : à $\gamma\in\SL_2(\Z)$ et à la courbe elliptique $E$ munie de
$\alpha:\Z^2\xrightarrow{\sim}R^1f_*\Z$, on associe $(E,\alpha\circ\gamma)$. Si l'on
regarde $X$, muni de
\[
q:\R^2\otimes_\R\cO_X\simeq R^1f_*\R\otimes_\R\cO_X\longrightarrow \omega^{-1},
\]
comme classifiant les structures complexes sur $\R^2$, on a la même action évidente du
groupe $\GL_2^+(\R)$ sur $(X,\R^2,\omega,q)$.

\paragraph{(2.3)}
Choisissons une base $(x_1,x_2)$ de $\R^2$ telle que $x_1\wedge x_2>0$. Un point
\[
(f:\R^2\to \C)\quad \bmod \C^\times
\]
de $X$ est repéré par
\[
z=\frac{f(x_1)}{f(x_2)}\qquad (\operatorname{Im}z>0),
\]
et l'application $q$ s'identifie à
\[
q:\R^2\to\Ga,\qquad ax_1+bx_2\longmapsto az+b.
\]
Ceci donne une trivialisation évidente de $\omega^{-1}$ sur $X$, qui n'est pas équivariante.
Relativement à cette trivialisation, une section $f(z)$ de $\omega^k$ sur $X$ est transformée
par un élément
\[
\begin{pmatrix}a&b\\ c&d\end{pmatrix}^{-1}\in \GL^+(\R^2)
\]
(relativement à la base $(x_1,x_2)$) suivant la formule
\[
f\cdot
\begin{pmatrix}a&b\\ c&d\end{pmatrix}^{-1}(z)
=(cz+d)^{-k}f\!\left(\frac{az+b}{cz+d}\right).
\]
De l'identité
\[
dz=(cz+d)^2\,d\!\left(\frac{az+b}{cz+d}\right)
\left|\begin{matrix}a&b\\ c&d\end{matrix}\right|^{-1},
\]
on déduit que $dz$ est une section de $\omega^{-2}\Omega^1_X$ invariante par
$\SL_2(\R)$. Cette section ne s'annule nulle part et définit un isomorphisme équivariant
de faisceaux $\SL_2(\R)$ entre $\omega^2$ et $\Omega^1_X$.

\paragraph{(2.4)}
Soit $\Gamma$ un sous-groupe discret de $\SL_2(\R)$, sans éléments d'ordre fini, et de
quotient de volume fini. On sait alors que l'espace quotient $X/\Gamma$ s'identifie à une
courbe projective lisse $\overline{X/\Gamma}$ privée d'un nombre fini de points. Le groupe
$\Gamma$ agit sur $X$ sans points fixes. Le système local équivariant $\R^2$ sur $X$,
ainsi que la suite exacte
\[
0\to \omega\to \R^2\otimes_\R\cO\to\omega^{-1}\to 0,
\]
définissent donc un système local $U$ sur $X/\Gamma$ et une suite exacte
\[
0\to \omega\to U\otimes_\R\cO\to\omega^{-1}\to 0.
\]

\paragraph{(2.5)}
Dans le cas particulier où $\Gamma\subset\SL_2(\Z)$, ces structures proviennent d'une
courbe elliptique $E$ sur $X/\Gamma$ dont l'image inverse est la courbe elliptique
équivariante $E_X$ sur $X$.

\paragraph{(2.6)}
Les points à l'infini de $X/\Gamma$ se décrivent comme suit (voir \([9]\)).

\begin{enumerate}[label=(\alph*)]
\item Ils correspondent aux classes de conjugaison dans $\Gamma$ de sous-groupes non
triviaux de $\Gamma$ maximaux parmi les sous-groupes de $\Gamma$ formés d'éléments
unipotents.

\item Soit $\Gamma_0\subset \Gamma$ un tel sous-groupe et choisissons une base
$(x_1,x_2)$ de $\R^2$ telle que, dans cette base, $\Gamma_0$ soit représenté par l'ensemble
de matrices
\[
\left\{\begin{pmatrix}1&0\\ n&1\end{pmatrix}\ \middle|\ n\in\Z\right\}.
\]
Soit $z$ la coordonnée (2.3) sur $X$ définie par $(x_1,x_2)$. Il existe $N$ tel que le
sous-ensemble
\[
X_N=\{z\mid \operatorname{Im}z>N\}
\]
de $X$ soit disjoint de tous ses conjugués $\gamma X_N$ pour $\gamma\in\Gamma/\Gamma_0$,
de sorte que $X_N/\Gamma_0\hookrightarrow X/\Gamma$. La fonction $q=e^{2\pi iz}$
établit un isomorphisme entre $X_N/\Gamma_0$ et le disque épointé
\[
0<|q|<e^{-2\pi N}.
\]
Si $P_{\Gamma_0}$ est le point de $\overline{X/\Gamma}-X/\Gamma$ associé à $\Gamma_0$,
cet isomorphisme se prolonge en un isomorphisme d'un voisinage de $P_{\Gamma_0}$ avec
le disque $0\le |q|<e^{-2\pi N}$.
\end{enumerate}

En vertu de (2.3), les sections de $\omega$ sur $X_N$ invariantes par $\Gamma_0$ s'identifient
aux fonctions holomorphes périodiques de période un sur $X_N$. On note encore $\omega$
le faisceau inversible sur $\overline{X/\Gamma}$ qui prolonge $\omega$ et tel que, au voisinage
d'un point $P_{\Gamma_0}$, la section de $\omega$ sur $X_N/\Gamma_0$ définie par la
fonction $1$ se prolonge en une section inversible sur $X_N/\Gamma_0$.

\paragraph{(2.7)}
Sur $X/\Gamma$ on a les deux faisceaux inversibles $\Omega^1$ et $\omega^2$, et un
isomorphisme $\varphi$ (2.3) entre leurs restrictions à $X/\Gamma$. De la formule
\[
dq=d(e^{2\pi iz})=2\pi i e^{2\pi iz}\,dz=2\pi iq\,dz
\]
il résulte que l'application
\[
\varphi:\Omega^1\to \omega^2
\]
se prolonge à $\overline{X/\Gamma}$, et présente un zéro simple en chacun des points à
l'infini.

\begin{definition}[2.8]
L'espace des formes paraboliques de poids $k+2$, relatives au groupe $\Gamma$, est l'espace
des sections globales
\[
H^0(\overline{X/\Gamma},\Omega^1\otimes\omega^k).
\]
En vertu de (2.7), cet espace s'identifie encore à l'espace des sections globales de
$\omega^{k+2}$ qui s'annulent à l'infini, c'est-à-dire en chaque « point ».
\end{definition}

\paragraph{(2.9)}
Notons $U^k$ la $k$-ième puissance symétrique du système local $U$ sur $X/\Gamma$.
L'application (2.5) induit une application
\[
\iota_k:\omega^k\to U^k\otimes_\R\cO,
\]
puis une application, encore notée $\iota_k$,
\[
\iota_k:\Omega^1\otimes\omega^k\to \Omega^1(U^k),
\]
où le but est le faisceau des formes différentielles holomorphes sur $X/\Gamma$, à
coefficients dans le système local $U^k$.

La résolution de de Rham de $U^k\otimes_\R\C$,
\[
0\to U^k\otimes_\R\C\to U^k\otimes_\R\cO
\xrightarrow{d}U^k\otimes_\R\Omega^1\to 0,
\]
induit une application
\[
\delta:H^0(X/\Gamma,\Omega^1(U^k))\to H^1(X/\Gamma,U^k\otimes\C).
\]
De plus, l'espace de cohomologie $H^1(X/\Gamma,U^k\otimes\C)$ est muni d'une
conjugaison complexe naturelle, telle que $\delta$ définisse une application antilinéaire
$\overline{\delta}$ depuis le conjugué complexe de $H^0(X/\Gamma,\Omega^1(U^k))$ vers
$H^1(X/\Gamma,U^k\otimes\C)$. On obtient ainsi une application
\[
\sh_0=\delta\circ H^0(\iota_k)\oplus
\overline{\delta}\circ H^0(\overline{\iota_k})
\]
\[
\sh_0:
H^0(\overline{X/\Gamma},\Omega^1\otimes\omega^k)
\oplus
\overline{H^0(\overline{X/\Gamma},\Omega^1\otimes\omega^k)}
\longrightarrow
H^1(X/\Gamma,U^k\otimes\C).
\]

Pour un faisceau arbitraire $F$ sur un espace $Y$, on note $\widetilde H^i(Y,F)$ l'image
de la cohomologie à supports compacts $H^i_c(Y,F)$ dans la cohomologie sans supports
$H^i(Y,F)$.

Le théorème 4.2.6 de \([9]\) est essentiellement équivalent au théorème suivant; dans
loc. cit. $k$ est supposé pair, mais la même démonstration vaut en général.

\begin{theorem}[2.10, Shimura {\rm [7]}]
Il existe un isomorphisme $\sh$ rendant commutatif le diagramme suivant :
\[
\begin{array}{ccc}
H^0(\overline{X/\Gamma},\Omega^1\otimes\omega^k)\oplus
\overline{H^0(\overline{X/\Gamma},\Omega^1\otimes\omega^k)}
& \xrightarrow{\ \sh\ } &
\widetilde H^1(X/\Gamma,U^k\otimes\C)
\\[0.4em]
\big\downarrow && \big\downarrow
\\[0.4em]
H^0(\overline{X/\Gamma},\Omega^1\otimes\omega^k)\oplus
\overline{H^0(\overline{X/\Gamma},\Omega^1\otimes\omega^k)}
& \xrightarrow{\ \sh_0\ } &
H^1(X/\Gamma,U^k\otimes\C).
\end{array}
\]
On appelle $\sh$ l'isomorphisme de Shimura.
\end{theorem}

\paragraph{(2.11)}
Dans le cas particulier où $\Gamma$ est un sous-groupe d'indice fini de $\SL_2(\Z)$, la
courbe elliptique $E$ sur $X/\Gamma$ fournit un schéma de courbe elliptique sur la courbe
algébrique $\overline{X/\Gamma}$; de manière équivalente, son invariant modulaire est
méromorphe à l'infini. Elle admet donc un modèle de Néron $E$ sur $\overline{X/\Gamma}$.
On peut montrer que les fibres de $E$ aux points à l'infini sont de type multiplicatif, et que
sur tout $\overline{X/\Gamma}$ on a
\[
\omega=e^*\Omega^1_{E/\overline{X/\Gamma}}.
\]
Dans ce cas particulier, on a $U=R^1f_*\Z\otimes\R$, de sorte que le but de l'isomorphisme
de Shimura peut s'écrire
\[
\widetilde H^1(X/\Gamma,U^k\otimes\C)
\simeq
\widetilde H^1(X/\Gamma,\Sym^k(R^1f_*\Z))\otimes_\Z\C .
\]

\section{Opérateurs de Hecke et représentation $\ell$-adique fondamentale}

\paragraph{(3.1)}
Rappelons (cf. \([3]\)) que la catégorie des faisceaux $\Z_\ell$ localement constants
constructibles, abrégés lcc, sur un schéma $S$ est la catégorie des systèmes projectifs
$(F_n)_{n\in\mathbb N}$ sur le site étale $S_{\et}$ de $S$ vérifiant :
\begin{enumerate}[label=(\roman*)]
\item $F_n$ est un faisceau localement constant de $\Z/(\ell^n)$-modules de type fini;
\item si $n\le m$, alors
\[
F_m\otimes_{\Z/(\ell^m)}\Z/(\ell^n)\xrightarrow{\sim}F_n.
\]
\end{enumerate}
Les faisceaux lcc en $\Z_\ell$ forment un champ de catégories abéliennes sur $S$; le
champ des faisceaux lcc en $\Q_\ell$ est le quotient de ce champ par le sous-champ épais
des faisceaux lcc en $\Z_\ell$ tués par une puissance de $\ell$. On note $\otimes\Q_\ell$
le foncteur canonique de la catégorie des faisceaux lcc en $\Z_\ell$ vers celle des faisceaux
lcc en $\Q_\ell$.

Si $S$ est connexe et muni d'un point géométrique $s$, la catégorie des faisceaux lcc en
$\Z_\ell$, respectivement en $\Q_\ell$, sur $S$ est équivalente, par le foncteur « fibre en
$s$ », à la catégorie des représentations continues du groupe fondamental $\pi_1(S,s)$ sur
des $\Z_\ell$-modules de type fini, respectivement sur des $\Q_\ell$-espaces vectoriels de
dimension finie.

Si $T$ est un ensemble fini de nombres premiers, un faisceau lcc en $\mathbb A_T^f$ consiste à se
donner, pour chaque nombre premier $\ell$, un faisceau lcc en $\Z_\ell$ si $\ell\notin T$,
et un faisceau lcc en $\Q_\ell$ si $\ell\in T$. Pour $T=\varnothing$, on parle de faisceaux
lcc en $\bZ$ plutôt que de faisceaux lcc en $\mathbb A_\varnothing^f$.

Pour $T$ général, la catégorie des faisceaux lcc en $\mathbb A_T^f$ est la limite inductive des
catégories de faisceaux lcc en $\mathbb A_{T'}^f$ pour $T'\subset T$ fini. On pose
\[
\Z_\ell=\varprojlim_n \Z/(\ell^n),\qquad
\Q_\ell=\Z_\ell\otimes\Q,\qquad
\bZ=(\Z_\ell)_\ell,
\]
et l'on utilise la notation adélique finie correspondante $\mathbb A_T^f$.

Le champ des courbes elliptiques à isogénie près sur $S$ est le champ obtenu à partir du
champ des courbes elliptiques sur $S$ en inversant formellement les isogénies. On note
$\otimes\Q$ le foncteur qui associe à une courbe elliptique la courbe elliptique sous-jacente
à isogénie près. Pour $S$ quasi-compact, on a
\[
\Hom(E,F)\otimes\Q\xrightarrow{\sim}\Hom(E\otimes\Q,F\otimes\Q),
\]
et, pour $S$ normal, toute courbe elliptique à isogénie près sur $S$ provient d'une courbe
elliptique sur $S$.

\paragraph{(3.2)}
Soit $f:E\to S$ une courbe elliptique sur un schéma $S$. On note $T_\ell(E)$ le système
projectif des noyaux $E_{\ell^n}$ de la multiplication par $\ell^n$ sur $E$, les applications
de transition de $E_{\ell^n}$ vers $E_{\ell^m}$, pour $n\ge m$, étant la multiplication par
$\ell^{\,n-m}$. En procédant de même pour $\Gm$, on pose $T_\ell(\Gm)=\Z_\ell(1)$.
Si $\ell$ est inversible sur $S$, alors $T_\ell(E)$ et $\Z_\ell(1)$ sont des faisceaux en
$\Z_\ell$ sur $S$. On définit $T_\infty(E)$ comme l'algèbre de Lie relative de $E$ sur
$S$, c'est-à-dire le dual inversible du faisceau inversible $\omega$ de (2.1(a)).

Supposons $S$ de caractéristique $0$. On définit alors le faisceau en $\bZ$, $T_f(E)$,
sur $S$ comme le système des $T_\ell(E)$, et l'on pose
\[
V_f(E)=T_f(E)\otimes \Af.
\]
Si $u:E\to F$ est une isogénie, alors $u$ induit un isomorphisme de $V_f(E)$ sur $V_f(F)$
et de $T_\infty(E)$ sur $T_\infty(F)$; les foncteurs $V_f$ et $T_\infty$ se factorisent donc
par la catégorie des courbes elliptiques à isogénie près sur $S$.

\begin{proposition}[3.3]
Soit $S$ un schéma de caractéristique $0$; soit $E_1(S)$ la catégorie des courbes elliptiques
sur $S$; et soit $E_2(S)$ la catégorie des triplets composés d'une courbe elliptique à
isogénie près $E$ sur $S$, d'un faisceau en $\bZ$, $T$, forme tordue de $\bZ^2$, et d'un
isomorphisme
\[
\beta:V_f(E)\xrightarrow{\sim}T\otimes\Af.
\]
Alors le foncteur
\[
I:E\longmapsto \bigl(E\otimes\Q,\ T_f(E),\ V_f(E)\simeq T_f(E)\otimes\Af\bigr)
\]
de $E_1(S)$ vers $E_2(S)$ est une équivalence de catégories.
\end{proposition}

La question est locale sur $S$, que l'on peut supposer quasi-compact. Si $f:E\to F$ est
un morphisme de courbes elliptiques sur $S$, et si $f$ est une isogénie, on a une suite
exacte
\[
0\to T_f(E)\to T_f(F)\to \ker(f)\to 0.
\tag{3.4}
\]
Le morphisme $f$ est divisible par $n$ si et seulement s'il tue le noyau $E_n$ de la
multiplication par $n$, parce que la multiplication par $n$ de $E/E_n$ vers $E$ est un
isomorphisme. D'après (3.4), ceci a lieu si et seulement si $T_f(f)$ est divisible par $n$,
et l'on en déduit que $\Hom_S(E,F)$ est le sous-groupe de $\Hom_S(E\otimes\Q,F\otimes\Q)$
formé des morphismes $f$ tels que $V_f(f)$ envoie $T_f(E)$ dans $T_f(F)$. Le foncteur
$I$ est donc pleinement fidèle.

Soit $X\in\operatorname{Ob}(E_2(S))$. Localement sur $S$, $X$ est défini par une courbe
elliptique à isogénie près $E\otimes\Q$, et par un réseau $T$ dans $V_f(E)$ qui, pour
presque tout $\ell$, coïncide avec $T_\ell(E)$. Pour $q\in\Q$, $(E\otimes\Q,T)$ est
isomorphe à $(E\otimes\Q,qT)$, ce qui permet de supposer $T_f(E)\subset T$. Le quotient
$K=T/T_f(E)$ est alors canoniquement isomorphe à un sous-groupe fini de $E$, et $X$ est
l'image de $E/K$ par $I$ (cf. 3.4). \QED

\begin{proposition}[3.5, Corollaire]
Le foncteur $F_1(S)$, respectivement $F'_1(S)$, qui associe à tout schéma $S$ de
caractéristique $0$ l'ensemble des classes d'isomorphisme de courbes elliptiques sur $S$
munies d'un isomorphisme
\[
\alpha:T_f(E)\xrightarrow{\sim}\bZ^2
\]
respectivement aussi d'un isomorphisme
\[
\alpha_\infty:T_\infty(E)\xrightarrow{\sim}\Ga,
\]
est isomorphe au foncteur $F_2(S)$, respectivement $F'_2(S)$, qui associe à $S$ l'ensemble
des classes d'isomorphisme de courbes elliptiques à isogénie près sur $S$, munies d'un
isomorphisme
\[
\beta:V_f(F)\xrightarrow{\sim}\Af^2
\]
respectivement aussi d'un isomorphisme
\[
\beta_\infty:T_\infty(F)\xrightarrow{\sim}\Ga.
\]
\end{proposition}

\begin{proposition}[3.6]
Le foncteur $F_1$ (respectivement $F'_1$) est représentable par un schéma $M_\infty$
(respectivement $M'_\infty$) sur $\Q$.
\end{proposition}

Soit $n\ge 3$ un entier. Le foncteur qui associe à chaque schéma $S$ l'ensemble des
classes d'isomorphisme de courbes elliptiques, munies d'un isomorphisme
\[
\alpha_n:E_n\xrightarrow{\sim}(\Z/n)^2
\]
respectivement aussi de $\alpha_\infty:T_\infty(E)\xrightarrow{\sim}\cO_X$, est représenté
par une courbe affine $M_n$ (respectivement par une surface affine $M'_n$) sur
$\Spec(\Z[1/n])$. Pour $n\mid m$, le morphisme de $M_m$ vers $M_n$ défini par
\[
(E,\alpha_m:E_m\xrightarrow{\sim}(\Z/m)^2)
\longmapsto
(E,(m/n)\alpha_m:E_n\xrightarrow{\sim}(\Z/n)^2)
\]
est fini et étale sur $\Spec(\Z[1/m])$, et l'on a
\[
M_\infty=\varprojlim M_n .
\]
On procède de même pour représenter $F'_1$. \QED

\paragraph{(3.7)}
Le schéma $M_\infty$ (respectivement $M'_\infty$) est muni d'une courbe elliptique
universelle $f_\infty:E_\infty\to M_\infty$ et d'un isomorphisme
\[
\alpha:T_f(E_\infty)\xrightarrow{\sim}\bZ^2,
\]
respectivement aussi d'un isomorphisme
\[
\alpha_\infty:T_\infty(E'_\infty)\xrightarrow{\sim}\Ga.
\]
D'après (3.5), le schéma $M_\infty$ (respectivement $M'_\infty$) représente le foncteur
$F_2$ (respectivement $F'_2$), ce qui donne une action évidente à gauche du groupe
adélique $\GL_2(\Af)$ sur
\[
(M_\infty,E_\infty\otimes\Q,\alpha\otimes\Af^2)
\]
respectivement sur
\[
(M'_\infty,E'_\infty,\alpha\otimes\Af^2,\alpha_\infty).
\]
Sur le foncteur, cette action est donnée, pour $g\in\GL_2(\Af)$, par
\[
g:(F,\beta:V_f(E)\xrightarrow{\sim}\Af^2,\beta_\infty)
\longmapsto
(F,g\circ\beta:V_f(E)\xrightarrow{\sim}\Af^2,\beta_\infty).
\]
Shafarevich a remarqué ce fait le premier.

Soit $Y$ un schéma sur $\C$, limite projective de schémas $Y_i$ de type fini sur $\C$,
les morphismes de transition étant finis. L'espace localement compact annelé $Y^{\an}$,
égal à la limite projective des $Y_i^{\an}$, ne dépend que de $Y$ et non de sa représentation
comme limite projective. Si $Y$ est un schéma sur $\Q$, limite projective de schémas $Y_i$
de type fini sur $\Q$ à morphismes de transition finis, on pose $Y^{\an}=(Y\otimes\C)^{\an}$.
Ceci s'applique à $M_\infty$ et à $M'_\infty$.

\begin{proposition}[3.8]
On a canoniquement
\[
M_\infty^{\prime\,\an}\simeq
\Hom(\Q^2\otimes\A,\C\times\Af^2)/\GL_2(\Q)
\]
et
\[
M_\infty^{\an}\simeq
\C^\times\backslash\Hom(\Q^2\otimes\A,\C\times\Af^2)/\GL_2(\Q),
\]
et, moins canoniquement,
\[
M_\infty^{\prime\,\an}\simeq \GL_2(\A)/\GL_2(\Q),\qquad
M_\infty^{\an}\simeq K_\infty\backslash \GL_2(\A)/\GL_2(\Q),
\]
où $K_\infty$ est le stabilisateur d'une structure complexe choisie à la place réelle. Ces
isomorphismes sont compatibles avec l'action de $\GL_2(\Af)$.
\end{proposition}

La notion de courbe elliptique à isogénie près, et ses variantes, s'étend au cas analytique
complexe. D'autre part, une isogénie $\varphi:E\to F$ induit un isomorphisme $\varphi^*$
entre les systèmes locaux de cohomologie rationnelle de $E$ et de $F$, ce qui permet de
définir ce dernier objet pour une courbe elliptique à isogénie près. Soit $S$ un espace
analytique complexe. À l'aide de (2.1), on voit que se donner une courbe elliptique à
isogénie près sur $S$ revient à se donner un faisceau inversible $T_\infty$, un système local
$T_\Q$ de $\Q$-espaces vectoriels, et un morphisme $\alpha:T_\Q\to T_\infty$ induisant
point par point un isomorphisme $T_\Q\otimes\R\to T_\infty$.

Soit $n$ un entier et soit $K_n$ le noyau de l'application naturelle
\[
\prod_\ell \GL_2(\Z_\ell)\to \GL_2(\Z/n).
\]
Soit $G_1$ le foncteur qui associe à $S$ l'ensemble des classes d'isomorphisme de courbes
elliptiques $f:E\to S$ sur $S$, munies d'un isomorphisme
\[
\varphi:\Q^2\xrightarrow{\sim}T_\Q(E),
\]
d'un isomorphisme $\alpha_\infty:T_\infty(E)\xrightarrow{\sim}\cO_S$, et d'un isomorphisme
$\alpha_n:E_n\xrightarrow{\sim}(\Z/n)^2$. Comme en (3.3), on voit que $G_1$ est
isomorphe au foncteur $G_2$ qui associe à $S$ l'ensemble des classes d'isomorphisme de
courbes elliptiques à isogénie près $E$ sur $S$, munies de
$\varphi:\Q^2\xrightarrow{\sim}T_\Q(E)$, de
$\alpha_\infty:T_\infty(E)\xrightarrow{\sim}\cO_S$, et d'un isomorphisme
$V_f(E)\xrightarrow{\sim}\Af^2$ donné localement sur $S$ à composition par un élément de
$K_n$ près.

Un tel objet est déterminé par l'application composée $\varphi'$ (donnée localement modulo
$K_n$) obtenue à partir de
\[
\varphi':\Q^2\xrightarrow{\varphi}T_\Q(E)
\longrightarrow T_\infty(E)\times V_f(E)
\xrightarrow{\sim}\cO_S\times \Af^2.
\]
On a
\[
E=\bZ^2\backslash \cO_S\times \Af^2/
\varphi'\bigl(\Q^2\cap \varphi'^{-1}(T_\infty(E)\times T_f(E))\bigr),
\]
de sorte que, cf. 2.2, $G_1$ et $G_2$ sont représentés par
\[
K_n\backslash\Isom(\Q^2\otimes\A,\C\times\Af^2).
\]

Supposons encore $n\ge 3$, de sorte que $\GL_2(\Q)$ agisse librement sur l'espace
précédent. L'espace analytique $M_n^{\an}$, respectivement $M_n^{\prime\,\an}$, représente
le foncteur analogue, en géométrie analytique, au foncteur représenté par $M_n$,
respectivement $M'_n$, parce que ce foncteur, disons $X$, est représentable et que la
flèche $X\to M_n^{\an}$, respectivement $X\to M_n^{\prime\,\an}$, induit une bijection sur
les ensembles de points à valeurs dans toute $\C$-algèbre de rang fini. On déduit donc de
ce qui précède que
\[
M_n^{\prime\,\an}\simeq
K_n\backslash\Isom(\Q^2\otimes\A,\C\times\Af^2)/\GL_2(\Q).
\]
En procédant de même pour $M_n$, on obtient la première assertion de (3.8) par passage
à la limite sur $n$.

Un point $x$ de
\[
\Isom(\Q^2\otimes\A,\C\times\Af^2)/\GL_2(\Q)
\]
s'identifie à un réseau $L_x$ de $\C\times\Af^2$, et la courbe correspondante est
\[
E_x\simeq \bZ^2\backslash \C\times\Af^2/L_x,
\]
munie de
\[
V_f(E)\simeq L_x\otimes\Af\xrightarrow{\sim}\Af^2.
\]
La dernière assertion de (3.8) en résulte facilement. \QED

Notons $f_n:E\to M_n$ la courbe elliptique universelle sur $M_n$. L'entier $k$ étant
fixé, on pose la définition suivante.

\begin{definition}[3.9]
On note $W$ (ou ${}_kW$ s'il peut y avoir confusion) le $\Q$-espace vectoriel
\[
W=\varinjlim_n\widetilde H^1(M_n^{\an},\Sym^k(R^1f_*\Q)).
\]
\end{definition}

Cet espace vectoriel ne dépend pas de la courbe elliptique universelle, à isogénie près,
$f_\infty:E\to M_\infty$; ainsi, par transport de structure, il est muni d'une action à gauche
du groupe adélique $\GL_2(\Af)$.

Si $\ell$ est un nombre premier, l'espace vectoriel $W_\ell=W\otimes\Q_\ell$ admet une
définition purement algébrique, en termes de cohomologie $\ell$-adique du schéma sur la
clôture algébrique $\overline\Q$ de $\Q$ obtenu à partir de $M_n$ par extension des scalaires :
\[
W_\ell=
\varinjlim_n
\widetilde H^1(M_n\otimes\overline\Q,\Sym^k(R^1f_{n*}\Q_\ell)).
\tag{3.10}
\]
Ainsi le groupe de Galois de $\overline\Q$ sur $\Q$ agit, par transport de structure, sur
$W_\ell$ et sur les espaces de niveau fini.

Enfin, l'espace $M_n^{\an}$ est la réunion disjointe de quotients du demi-plan de Poincaré
par des sous-groupes de congruence de $\SL_2(\Z)$; donc, en notant $\omega$ le faisceau
inversible sur $M_n$ défini par $E$, la théorie de Shimura (2.10) donne
\[
W_\infty=W\otimes\C=
\varinjlim_n\left(
H^0(M_n^{\an},\Omega^1\otimes\omega^k)
\oplus
\overline{H^0(M_n^{\an},\Omega^1\otimes\omega^k)}
\right).
\tag{3.11}
\]
Cette décomposition de $W\otimes\C$ en deux sous-espaces conjugués complexes, dont l'un
est l'espace de toutes les formes paraboliques de poids $k+2$ relatives à un sous-groupe de
congruence général de $\SL_2(\Z)$, est analogue à une décomposition de Hodge, de type
\[
(0,k+1)+(k+1,0).
\]
L'action du groupe adélique commute avec celle du groupe de Galois et respecte la
décomposition précédente.

Bien que le système local $\ell$-adique $R^1f_*\Q_\ell$ soit trivial sur $M_\infty$, je ne
sais pas s'il existe une relation entre $W_\ell$ et
\[
\varinjlim_n\bigl(\widetilde H^1(M_n\otimes\overline\Q,\Q_\ell)\bigr)
\otimes \Sym^k(\Q_\ell^2).
\]

\paragraph{(3.12)}
Soit $n\ge 3$ un entier et soit $K_n$ comme en (3.8). On a alors
\[
W^{K_n}={}_nW.
\]
Ceci se vérifie par passage à la limite, et résulte du fait qu'en cohomologie rationnelle, la
cohomologie du quotient d'un espace par un groupe fini s'obtient en prenant les invariants
de ce groupe dans la cohomologie.

Posons, pour $p$ premier,
\[
W^{(p)}=W^{\GL_2(\Z_p)}.
\]
Par passage à la limite, on obtient
\[
W^{(p)}=\varinjlim_{(n,p)=1}{}_nW.
\]
Sur cet espace de cohomologie agissent encore :
\begin{enumerate}[label=(\roman*)]
\item le sous-groupe $\prod_{\ell\ne p}\GL_2(\Q_\ell)$ de $\GL_2(\Af)$, parce que ce
sous-groupe centralise $\GL_2(\Z_p)$;
\item l'algèbre de Hecke $\cH(\GL_2(\Q_p),\GL_2(\Z_p))$, qui est l'algèbre des mesures
entières sur l'espace discret $\GL_2(\Q_p)/\GL_2(\Z_p)$ invariantes à gauche par
$\GL_2(\Z_p)$. Cette sous-algèbre de l'algèbre de groupe de $\GL_2(\Q_p)$ agit sur
$W$ et respecte $W^{(p)}$. Elle agit déjà sur chaque ${}_nW$ pour $n$ premier à $p$.
\end{enumerate}

L'algèbre de Hecke admet pour base les mesures associées aux fonctions caractéristiques
des doubles classes de $\GL_2(\Z_p)$ dans $\GL_2(\Q_p)$, et l'on sait que
\[
\cH(\GL_2(\Q_p),\GL_2(\Z_p))=\Z[T_p,R_p,R_p^{-1}],
\]
où $T_p$ et $R_p$ sont les doubles classes de
\[
\begin{pmatrix}1&0\\0&p^{-1}\end{pmatrix}
\qquad\text{et}\qquad
\begin{pmatrix}p^{-1}&0\\0&p^{-1}\end{pmatrix}.
\]


\paragraph{(3.13)}
Soient $p$ un nombre premier, $n\geq 3$ un entier premier à $p$, et $F_{n,p}$ le foncteur qui associe à tout schéma $S$ l'ensemble des classes d'isomorphisme de diagrammes commutatifs de $S$-schémas
\[
\begin{array}{ccccc}
&& \Z/(n)^2 &&\\[-0.2em]
& \alpha\nearrow && \nwarrow\alpha' &\\[-0.2em]
E_n & \longrightarrow & F_n &&\\[-0.2em]
\cap && \cap &&\\[-0.2em]
E & \xrightarrow{\ \varphi\ } & F &&
\end{array}
\tag{3.14}
\]
où $\varphi$ est une $p$-isogénie entre courbes elliptiques et $\alpha$ un isomorphisme. On note $q_1$ et $q_2:F_{n,p}\to M_n$ les morphismes de foncteurs qui associent à un diagramme (3.14) les sous-diagrammes $(E,E_n,\alpha)$ et $(F,F_n,\alpha')$.

\begin{proposition}[3.15]
Le foncteur $F_{n,p}$ est représenté par un schéma $M_{n,p}$, et les morphismes $q_1,q_2:M_{n,p}\to M_n$ sont finis.
\end{proposition}

L'automorphisme $\sigma$ de $F_{n,p}$ qui envoie $\varphi:E\to F$ sur ${}^{t}\!\varphi:F\to E$ échange $q_1$ et $q_2$; il suffit donc de considérer $q_1$. Ce morphisme identifie $F_{n,p}$ au foncteur des sous-groupes d'ordre $p$ de la courbe elliptique universelle $E$ sur $M_n$; ainsi, par la théorie des schémas de Hilbert, $F_{n,p}$ est représentable et $M_{n,p}$ est propre sur $M_n$. Si $s$ est un point géométrique de $M_n$, $q_1^{-1}(s)$ est l'ensemble des sous-groupes d'ordre $p$ de $E_s$; il a $p+1$ éléments si $\charop(k(s))\ne p$, et un seul, le noyau de Frobenius, si $\charop(k(s))=p$. \QED

On peut montrer que $M_{n,p}$ est régulier, et que $q_1$ et $q_2$ sont finis plats; nous n'utiliserons pas ce résultat délicat, mais nous nous contentons de noter que, sur $\Spec(\Z[1/p])$, chacun des $q_i$ fait de $M_{n,p}$ un revêtement étale de degré $p+1$ de $M_n$.

Les morphismes $q_i$ s'insèrent dans un diagramme commutatif
\[
\begin{array}{ccccc}
q_1^*E & \xrightarrow{\ \varphi\ } & q_2^*E &&\\[-0.2em]
\swarrow && \searrow &&\\[-0.2em]
E & \xleftarrow{\ u\ } & M_{n,p} & \xrightarrow{\ v\ } & E\\[-0.2em]
\downarrow f_n && \downarrow q_1\quad\downarrow q_2 && \downarrow f_n\\[-0.2em]
M_n && M_n && M_n
\end{array}
\tag{3.16}
\]
où $(\varphi,u,v)$ fait partie du diagramme universel (3.14).

On note $I_p$ le morphisme de $M_n$ vers $M_n$ correspondant au morphisme de foncteurs $(E,\alpha)\mapsto(E,\alpha/p)$:
\[
\begin{array}{ccc}
E&\longrightarrow&E\\
\downarrow&&\downarrow\\
M_n&\xrightarrow{\ I_p\ }&M_n
\end{array}
\qquad
I_p^*(E,\alpha)=(E,\alpha/p).
\tag{3.17}
\]
$I_p^*$ est un automorphisme de $\widetilde H^i(M_n^{\an},\Sym^k(R^1f_{n*}\Z))$. Il est fastidieux, mais routinier, de prouver ce qui suit.

\begin{proposition}[3.18]
\leavevmode
\begin{enumerate}[label=(\roman*)]
\item L'endomorphisme $T_p$ de ${}_nW$ s'exprime, à l'aide de (3.16), comme l'application composée
\[
\begin{aligned}
&\widetilde H^1(M_n^{\an},\Sym^k(R^1f_{n*}\Q))
\xrightarrow{\ q_2^*\ }
\widetilde H^1(M_{n,p}^{\an},\Sym^k(R^1v_*\Q))\\
&\xrightarrow{\ \varphi^*\ }
\widetilde H^1(M_{n,p}^{\an},\Sym^k(R^1u_*\Q))
\xrightarrow{\ q_{1*}\ }
\widetilde H^1(M_n^{\an},\Sym^k(R^1f_{n*}\Q)).
\end{aligned}
\]
où $q_{1*}$ est le morphisme trace du revêtement $q_1$.
\item De même, $R_p=p^kI_p^*$.
\end{enumerate}
\end{proposition}
\QED

Le lecteur méfiant pourrait oublier les préliminaires adéliques et définir $T_p$ par (i). Pour $n=1$ ou $2$, on pose ${}_nW=W^{K_n}$, de sorte que
\[
{}_1W={}_nW^{\GL_2(\Z/(n))}.
\]
Si $S_{k+2}$ désigne l'espace des formes paraboliques, pour le groupe $\SL_2(\Z)$, de poids $k+2$, l'isomorphisme de Shimura (3.11) induit un isomorphisme
\[
{}^k_1W_\infty={}^k_1W\otimes\C=S_{k+2}\oplus \overline{S}_{k+2}.
\]
Il est fastidieux, mais routinier, de prouver ce qui suit.

\begin{proposition}[3.19]
L'endomorphisme $T_p$ de ${}^k_1W_\infty$ s'identifie, par l'isomorphisme de Shimura, à la somme directe de l'opérateur de Hecke sur $S_{k+2}$, avec le facteur $p^{k-1}$, et de son conjugué.
\end{proposition}
\QED

\paragraph{(3.20)}
On a canoniquement
\[
\bigwedge^2R^1f_{n*}\Z_\ell\simeq R^2f_{n*}\Z_\ell\simeq \Z_\ell(-1),
\]
de sorte que $\Sym^k(R^1f_{n*}\Z_\ell)$ est muni d'une forme bilinéaire, symétrique pour $k$ pair et alternée pour $k$ impair, à valeurs dans $\Z_\ell(-k)$. La forme obtenue par tensorisation avec $\Q_\ell$ est non dégénérée.

Si $F$ est un faisceau lcc en $\Q_\ell$ sur un schéma lisse $X$, de dimension pure $n$ sur un corps algébriquement clos $k$, la dualité de Poincaré donne
\[
H^i(X,F)^\vee\simeq H_c^{2n-i}(X,\Hom(F,\Q_\ell(n))),
\]
\[
H_c^i(X,F)^\vee\simeq H^{2n-i}(X,\Hom(F,\Q_\ell(n))),
\]
\[
\widetilde H^i(X,F)^\vee\simeq \widetilde H^{2n-i}(X,\Hom(F,\Q_\ell(n))).
\]
En prenant $X=\overline M_n$ et $F=\Sym^k(R^1f_{n*}\Q_\ell)$, on définit une forme bilinéaire non dégénérée ${}_n(\ ,\ )$ sur ${}^k_nW_\ell$, à valeurs dans $\Q_\ell(-k-1)$. Cette forme est symétrique pour $k$ impair et alternée pour $k$ pair. C'est l'analogue $\ell$-adique du produit intérieur de Peterson. Si $n\mid m$, et si le revêtement $\psi:M_m\to M_n$ est de degré $d$, alors
\[
{}_m(\psi^*x,\psi^*y)=d\cdot {}_n(x,y).
\]

\section{La formule de congruence}

Dans cette section, on fixe des entiers $k\geq 0$ et $n\geq 3$, et des nombres premiers $p$ et $\ell$. On suppose que $p$ est premier à $\ell$ et à $n$. On note $f_n:E\to M_n$ la courbe elliptique universelle sur $M_n$, munie de $\alpha:E_n\xrightarrow{\sim}\Z/(n)^2$.

Quel que soit le schéma $Y$, on note $a$ l'unique morphisme de $Y$ vers $\Spec(\Z)$, ou, si nécessaire, vers un sous-schéma de $\Spec(\Z)$. Si $Y$ est séparé et de type fini sur $\Spec(\Z)$, et si $F$ est un faisceau en $\Z_\ell$ ou en $\Q_\ell$ sur $Y$, on note $R^ia_*(Y,F)$, respectivement $R^ia_!(Y,F)$, respectivement $R^i\widetilde a(Y,F)$, le faisceau en $\Z_\ell$ ou en $\Q_\ell$ sur $\Spec(\Z)$ qui est la $i$-ième image directe supérieure de $F$ par $a$, respectivement la $i$-ième image directe à supports propres, respectivement $\Imop(R^ia_!(Y,F)\to R^ia_*(Y,F))$. Pour $m\in\mathbb N$, on pose
\[
Y[1/m]=Y\times\Spec(\Z[1/m]).
\]

\begin{theorem}[4.1, Igusa {\rm [1]}]
Le schéma $M_n$ peut être compactifié en un schéma de courbes $M_n^*$, projectif et lisse sur $\Spec(\Z[1/n])$, tel que $M_n^*\setminus M_n$ soit un revêtement étale de $\Spec(\Z[1/n])$.
\end{theorem}

Les schémas $M_n$ sont formellement lisses, donc lisses sur $\Spec(\Z)$. L'invariant modulaire $j$ de la courbe universelle sur $M_n$ est un morphisme de $M_n$ vers la droite affine $\mathbb A^1$ sur $\Spec(\Z[1/n])$. Le morphisme $j$ est fini et est un revêtement étale hors des sections $0$ et $1728$ de $\mathbb A^1$; en effet:
\begin{enumerate}[label=(\alph*)]
\item Deux courbes elliptiques sur un corps algébriquement clos qui ont le même invariant $j$ sont isomorphes (par exemple [8] 6.3); les fibres géométriques de $j$ sont donc finies. Les schémas $M_n$ et $\mathbb A^1$ étant lisses de même dimension sur $\Spec(\Z)$, $j$ est quasi-fini et plat.
\item Si $E$ est une courbe elliptique sur le corps des fractions $K$ d'un anneau de valuation discrète $R$, dont l'invariant $j$ appartient à $R$, et dont les points d'ordre $n$ sont rationnels sur $K$, alors $E$ a bonne réduction. Le critère valuatif de propreté montre donc que $j$ est propre.
\item Si $E$ et $F$ sont deux courbes elliptiques sur un schéma $S$, de même invariant $j$, et si $j$ et $j-1728$ sont inversibles, alors le schéma $\Isom(S;E,F)$ des isomorphismes entre $E$ et $F$ est étale sur $S$ ([8] 6.3). Dans le diagramme
\[
\begin{array}{ccc}
\Isom(M_n\times_{\mathbb A^1}M_n;\operatorname{pr}_1^*E,\operatorname{pr}_2^*E) & \sim & M_n\times\GL_2(\Z/(n))\\
\downarrow u && \downarrow v\\
M_n\times_{\mathbb A^1}M_n&\xrightarrow{\operatorname{pr}_1}&M_n,
\end{array}
\]
où $j\ne0,1728$, $u$ et $v$ sont étales surjectifs; donc $\operatorname{pr}_1$ est étale et, par descente fidèlement plate, $j$ est étale.
\end{enumerate}
La section à l'infini de la droite projective $\mathbb P^1\supset\mathbb A^1$ sur $\Spec(\Z[1/n])$ est un diviseur régulier, de point générique de caractéristique $0$, dans un schéma régulier. Il résulte alors d'un théorème d'Abyankhar (voir [5]) que, le long de ce diviseur $j=\infty$, $M_n$ est modérément, ou docilement, ramifié au-dessus de $\mathbb P^1$, et que la normalisation $M_n^*$ de $\mathbb P^1$ dans $M_n$ vérifie (4.1). \QED

Il résulte du même théorème que les faisceaux lcc en $\Z_\ell$, $R^if_{n*}\Z_\ell$, sur $M_n[1/\ell]$ sont modérément ramifiés à l'infini. Par conséquent, d'après (4.1) et les théorèmes de spécialisation pour la cohomologie $\ell$-adique (voir [5]), $R^ia_*(M_n,\Sym^k(R^1f_{n*}\Z_\ell))$, $R^ia_!(M_n,\Sym^k(R^1f_{n*}\Z_\ell))$, et donc $R^i\widetilde a(M_n,\Sym^k(R^1f_{n*}\Z_\ell))$ sont des faisceaux lcc en $\Z_\ell$ sur $\Spec(\Z[1/n,1/\ell])$, dont la formation est compatible avec tout changement de base.

\begin{proposition}[4.2, Corollaire]
Le module galoisien ${}_nW_\ell$ est la fibre du faisceau lcc en $\Q_\ell$
\[
R^1\widetilde a(M_n,\Sym^k(R^1f_{n*}\Z_\ell))\otimes\Q_\ell
\]
sur $\Spec(\Z[1/n,1/\ell])$ au point géométrique $\overline\Q$. Il est non ramifié hors de $n$ et $\ell$.
\end{proposition}
\QED

Écrivons les deux diagrammes commutatifs sur $M_n\otimes\F_p$
\[
\begin{array}{ccccc}
&&\Z/(n)^2&&\\
&\alpha\nearrow&&\nwarrow\alpha^{(p)}&\\
E_n&\xrightarrow{\ F\ }&E_n^{(p)}&&\\
\cap&&\cap&&\\
E&\xrightarrow{\ F\ }&E^{(p)}&&
\end{array}
\qquad
\begin{array}{ccccc}
&&\Z/(n)^2&&\\
&p\alpha^{(p)}\nearrow&&\nwarrow\alpha&\\
E_n^{(p)}&\xrightarrow{\ V\ }&E_n&&\\
\cap&&\cap&&\\
E^{(p)}&\xrightarrow{\ V\ }&E&&
\end{array}
\]
plus brièvement sous la forme
\[
F:(E,\alpha)\to(E^{(p)},\alpha^{(p)})
\qquad\text{et}\qquad
V:(E^{(p)},p\alpha^{(p)})\to(E,\alpha),
\]
où $F$ est le morphisme de Frobenius et $V$, son transposé, est le Verschiebung. Ces diagrammes définissent des morphismes $\Phi_1$ et $\Phi_2$ de $M_n\otimes\F_p$ vers $M_{n,p}$. Ces morphismes sont finis, considérés comme sections de $q_1$ et $q_2$, et définissent un morphisme
\[
\Phi=\Phi_1\coprod\Phi_2:M_n\otimes\F_p\coprod M_n\otimes\F_p\to M_{n,p}\otimes\F_p.
\]
Soit $\Phi^h$ la restriction de $\Phi$ aux ouverts $M_n^h$ et $M_{n,p}^h$ de $M_n\otimes\F_p$ et de $M_{n,p}\otimes\F_p$ correspondant aux courbes d'invariant de Hasse $h$ non nul.

\begin{proposition}[4.3]
$\Phi^h$ est un isomorphisme.
\end{proposition}

Soit $\varphi:E_1\to E_2$ une $p$-isogénie entre courbes elliptiques d'invariant de Hasse inversible sur un schéma $S$ de caractéristique $p$. En chaque point géométrique de $S$, ou bien le noyau $\ker(\varphi)$ de $\varphi$ est étale sur $S$, ou bien son dual de Cartier, isomorphe à $\ker({}^{t}\!\varphi)$, est étale sur $S$. La propriété « $\ker(\varphi)$ est étale » est ouverte; ainsi, localement sur $S$, ou bien $\ker(\varphi)$ est purement infinitésimal, ou bien $\ker({}^{t}\!\varphi)$ l'est. Le seul sous-groupe infinitésimal d'ordre $p$ de $E_1$ ou de $E_2$ étant le noyau de Frobenius, dans le premier cas $\varphi$ est isomorphe à $F:E_1\to E_1^{(p)}$, et dans le second cas ${}^{t}\!\varphi$ est isomorphe à $F:E_2\to E_2^{(p)}$ et $\varphi$ à $V:E_2^{(p)}\to E_2$. \QED

\begin{proposition}[4.4]
\leavevmode
\begin{enumerate}[label=(\roman*)]
\item Le schéma $M_{n,p}$ est lisse sur $\Spec(\Z)$ hors des points de caractéristique $p$ où $h=0$.
\item Les morphismes $q_1$ et $q_2$ induisent des morphismes finis plats $q'_1$ et $q'_2$ de la normalisation $M'_{n,p}$ de $M_{n,p}$ vers $M_n$.
\item Le morphisme $\Phi$ se factorise par un morphisme surjectif
\[
\Phi':M_n\otimes\F_p\coprod M_n\otimes\F_p\to M'_{n,p}\otimes\F_p.
\]
\end{enumerate}
\end{proposition}

L'automorphisme $\sigma$ de (3.15) échange $\varphi$ et ${}^{t}\!\varphi$; il suffit donc de prouver (i) aux points de caractéristique $p$ de $M_{n,p}$ où le noyau de $\varphi$ est infinitésimal: il n'y a pas d'obstruction au relèvement infinitésimal d'une courbe elliptique et de la partie infinitésimale du noyau de la multiplication par $p$.

Là où $p=h=0$, la fibre du morphisme fini (3.15) $q_i:M_{n,p}\to M_n$ est réduite à un point; l'ouvert de lissité de $M_{n,p}$ est donc dense dans $M_{n,p}$, et $M'_{n,p}$ est partout de dimension $2$. Le schéma $M_n$ étant régulier, d'après EGA $0_{\mathrm{IV}}$ 16.5.1 et 17.3.5(ii), le morphisme $q_i:M'_{n,p}\to M_n$ est plat. Enfin, (iii) résulte du fait que $\Phi$ est fini et que $M_n\otimes\F_p$ est une courbe normale. \QED

L'endomorphisme de Hecke $T_p$ de ${}_nW_\ell$, explicité en (3.18), est la tensorisation par $\Q_\ell$ de la fibre au point géométrique $\overline\Q$ de $\Spec(\Z[1/n,1/\ell])$ de l'endomorphisme, encore noté $T_p$, de $R^1\widetilde a(M_n,\Sym^k(R^1f_{n*}\Z_\ell))$ défini par la correspondance
\[
\begin{array}{ccccc}
q_1'^*E&\xrightarrow{\ \varphi\ }&q_2'^*E&&\\
\swarrow&&\searrow&&\\
E&\xleftarrow{\ u\ }&M'_{n,p}&\xrightarrow{\ v\ }&E\\
\downarrow f_n&&\downarrow q'_1\quad\downarrow q'_2&&\downarrow f_n\\
M_n&&M_n&&M_n,
\end{array}
\tag{4.5}
\]
avec
\[
T_p=q'_{1*}\,\varphi^*\,q_2'^* \qquad (\text{cf. }3.18).
\]
Les endomorphismes $R_p$ et $I_p$ s'interprètent de même.

\begin{lemma}[4.6]
Soient $X,Y,Z_1$ et $Z_2$ quatre schémas séparés et de type fini sur un schéma noethérien $S$; soit $F$ un faisceau en $\Z_\ell$ sur $X$, soit $G$ un faisceau en $\Z_\ell$ sur $Y$, et notons $a$ chacun des morphismes structuraux de $X,Y,Z_1$ ou $Z_2$ vers $S$. Supposons donné le diagramme commutatif suivant de $S$-schémas et de morphismes de faisceaux:
\[
\begin{array}{ccccc}
y_1^*G&\xrightarrow{z_1}&x_1^*F&&y_2^*G\xrightarrow{z_2}x_2^*F\\
&&&&\\[-1em]
Z_1&\xrightarrow{f}&Z_2&&Z_1\xrightarrow{f}Z_2\\[-0.2em]
\downarrow x_1&&\downarrow x_2&&\downarrow y_1\quad\downarrow y_2\\[-0.2em]
X&&X&&Y
\end{array}
\]
Supposons que $f^*z_2=z_1$, que $y_1$ et $y_2$ soient propres, que $x_1$ et $x_2$ soient finis et plats, et que, pour tout point géométrique $s$ de $Z_2$, la multiplicité de $s$ dans sa fibre $x_2^{-1}(x_2(s))$ soit égale à la somme des multiplicités, dans leurs fibres pour $x_1$, des points géométriques de $Z_1$ dans l'image inverse de $s$ par $f$. Alors le diagramme
\[
\begin{array}{ccccc}
R^i\widetilde a(Y,G)&\xrightarrow{y_1^*}&R^i\widetilde a(Z_1,G)&\xrightarrow{z_1}&R^i\widetilde a(Z_1,F)\xrightarrow{x_{1*}}R^i\widetilde a(X,F)\\
\Vert&&\uparrow f^*&&\uparrow f^*\qquad\Vert\\
R^i\widetilde a(Y,G)&\xrightarrow{y_2^*}&R^i\widetilde a(Z_2,G)&\xrightarrow{z_2}&R^i\widetilde a(Z_2,F)\xrightarrow{x_{2*}}R^i\widetilde a(X,F)
\end{array}
\]
est commutatif.
\end{lemma}

Ce lemme résulte des lemmes analogues pour $R^ia_!$ et $R^ia_*$. La commutativité des premiers carrés est triviale. Le dernier se récrit
\[
\begin{array}{ccc}
R^i\widetilde a(Z_1,x_1^*F)&\xleftarrow{\sim}&R^i\widetilde a(X,x_{1*}x_1^*F)\xrightarrow{\Tr}R^i\widetilde a(X,F)\\
\uparrow&&\uparrow\qquad\Vert\\
R^i\widetilde a(Z_2,x_2^*F)&\xleftarrow{\sim}&R^i\widetilde a(X,x_{2*}x_2^*F)\xrightarrow{\Tr}R^i\widetilde a(X,F),
\end{array}
\]
et l'on rappelle la définition de la trace pour vérifier que le carré
\[
\begin{array}{ccc}
x_{1*}x_1^*F&\xrightarrow{\Tr}&F\\
\uparrow&&\Vert\\
x_{2*}x_2^*F&\xrightarrow{\Tr}&F
\end{array}
\]
est commutatif. \QED

\paragraph{(4.7)}
On note $T_p/\F_p$ l'endomorphisme induit par $T_p$ sur la restriction à $\Spec(\F_p)$ du faisceau lcc en $\Z_\ell$ $R^1\widetilde a(M_n,\Sym^kR^1f_{n*}\Z_\ell)$. On a
\[
R^1\widetilde a(M_n,\Sym^kR^1f_{n*}\Z_\ell)|\Spec(\F_p)
\xrightarrow{\sim}
R^1\widetilde a(M_n\otimes\F_p,\Sym^kR^1f_{n*}\Z_\ell);
\]
la formation du morphisme trace pour un morphisme fini plat est compatible au changement de base. Ainsi $T_p/\F_p$ peut être construit, sur le modèle (3.18), à partir de la fibre sur $\F_p$ de la correspondance (4.5). Le lemme (4.6), appliqué aux diagrammes commutatifs
\[
\begin{array}{ccc}
M_n\otimes\F_p\coprod M_n\otimes\F_p&\xrightarrow{\Phi'}&M'_{n,p}\otimes\F_p\\
\downarrow&&\swarrow q'_1\\
M_n\otimes\F_p&&
\end{array}
\qquad
\begin{array}{ccc}
M_n\otimes\F_p\coprod M_n\otimes\F_p&\xrightarrow{\Phi'}&M'_{n,p}\otimes\F_p\\
&&q'_2\searrow\quad\downarrow\\
&&M_n\otimes\F_p,
\end{array}
\]
fournit alors une décomposition de $T_p/\F_p$ en somme d'endomorphismes définis par les deux correspondances suivantes:
\[
\begin{array}{cccccc}
&& (E,\alpha)&\xrightarrow{F}&(E^{(p)},\alpha^{(p)})=F^*(E,\alpha)\\[-0.2em]
(a)&& \swarrow && \searrow\\[-0.2em]
& (E,\alpha) && M_n\otimes\F_p && (E,\alpha)\\[-0.2em]
&& \searrow && \swarrow F\\[-0.2em]
&& M_n\otimes\F_p && M_n\otimes\F_p
\end{array}
\]
où $F$ est le Frobenius absolu. On reconnaît dans cette correspondance le Frobenius géométrique
\[
\begin{array}{cccccc}
&&x^*(E,\alpha)=(E^{(p)},p\alpha^{(p)})&\xrightarrow{\ V\ }&(E,\alpha)\\[-0.2em]
(b)&& \swarrow && \searrow\\[-0.2em]
& (E,\alpha) && M_n\otimes\F_p && (E,\alpha)\\[-0.2em]
&& \searrow x && \swarrow\\[-0.2em]
&& M_n\otimes\F_p && M_n\otimes\F_p.
\end{array}
\]
L'application $x$ est la composée $I_p^{-1}\circ F$:
\[
\begin{array}{ccccc}
(E,\alpha)&\longleftarrow&(E,p\alpha)&\longleftarrow&(E^{(p)},p\alpha^{(p)})\\
\downarrow&&\downarrow&&\downarrow\\
M_n\otimes\F_p&\xleftarrow{\ I_p^{-1}\ }&M_n\otimes\F_p&\xleftarrow{\ F\ }&M_n\otimes\F_p.
\end{array}
\]
L'endomorphisme correspondant est donc le composé de
\[
V:
R^1\widetilde a(M_n\otimes\F_p,\Sym^k(R^1f_{n*}\Z_\ell))
\xrightarrow{V^*}
R^1\widetilde a(M_n\otimes\F_p,\Sym^k(R^1f^{(p)}_{n*}\Z_\ell))
\xrightarrow{\Tr_F}
R^1\widetilde a(M_n\otimes\F_p,\Sym^k(R^1f_{n*}\Z_\ell))
\]
et de
\[
I_p^*=\Tr_{I_p^{-1}}:\quad\text{endomorphisme de }R^1\widetilde a(M_n\otimes\F_p,\Sym^kR^1f_{n*}\Z_\ell).
\]

\begin{proposition}[4.8]
On a $T_p/\F_p=F+I_p^*V$, et:
\begin{enumerate}[label=(\roman*)]
\item $F$ s'identifie à l'inverse de l'élément de Frobenius arithmétique $\varphi_p$ du groupe de Galois $\Gal(\overline\F_p/\F_p)$ agissant sur $\widetilde H^1(M_n\otimes\overline\F_p,\Sym^k(R^1f_{n*}\Z_\ell))$;
\item $F$ et $V$ sont transposés l'un de l'autre relativement au produit scalaire (3.20);
\item $FV=VF=p^{k+1}$.
\end{enumerate}
\end{proposition}

Pour la relation (i) entre le Frobenius géométrique et le Frobenius arithmétique, on renvoie à l'exposé de C. Houzel (SGA 5.XV). Le composé $VF$ est le composé des homomorphismes obtenus des applications suivantes:
\[
\begin{array}{ccccc}
E&\xleftarrow{F_E}&E^{(p)}&\xrightarrow{V_E}&E\\
\downarrow&&\downarrow&&\downarrow\\
M_n&\xrightarrow{F}&M_n&\xrightarrow{F}&M_n.
\end{array}
\]
Ainsi
\[
VF=\Tr_F\circ F_E^*\circ V_E^*\circ F^*.
\]
L'application $F_E^*V_E^*=(F_EV_E)^*=(p\cdot1_E)^*$ agit par multiplication par $p^k$ sur $\Sym^kR^1f_{n*}\Z_\ell$; ainsi
\[
VF=p^k\Tr_F\circ F^*=p^k\cdot p=p^{k+1},
\]
car $F:M_n\to M_n$ est de degré $p$.

Par transport de structure, $\varphi_p$ respecte le produit scalaire (3.20), à valeurs dans $\Q_\ell(-k-1)$, groupe sur lequel $\varphi_p$ agit par multiplication par $p^{-k-1}$. On a donc
\[
(Fx,y)=p^{k+1}(\varphi_pFx,\varphi_py)=(x,p^{k+1}F^{-1}y)=(x,Vy).
\]
\QED

Le théorème suivant, équivalent à (4.8), remonte à Eichler.

\begin{theorem}[4.9, formule de congruence]
Soit $K_{n,\ell}$ la plus grande sous-extension de $\overline\Q$ non ramifiée hors de $n$ et $\ell$; soit $\varphi_p$ un élément de Frobenius relatif à $p$ dans $\Gal(K_{n,\ell}/\Q)$; soit $F$ l'endomorphisme $\varphi_p^{-1}$ de ${}_nW_\ell$; et soit $V$ le transposé de $F$ relativement au produit scalaire (3.20). Alors
\[
T_p=F+I_p^*V,
\qquad
FV=p^{k+1},
\qquad
1-T_pX+pR_pX^2=(1-FX)(1-I_p^*VX).
\]
\end{theorem}
\QED

\section{Weil implique Ramanujan}

Si $p$ est un nombre premier et $X$ un schéma sur $\F_p$, on note $\overline\F_p$ une clôture algébrique de $\F_p$, $F:X\to X$ l'endomorphisme de Frobenius géométrique, et l'on pose $\overline X=X\otimes\overline\F_p$; $\ell$ désigne toujours un nombre premier différent de $p$.

Par les « conjectures de Weil », on entend l'énoncé suivant: si $X$ est un schéma projectif lisse sur $\F_p$ et si $\ell$ est un nombre premier différent de $p$, alors les valeurs propres de l'endomorphisme $F^*$ de $H^i(\overline X,\Q_\ell)$ sont des entiers algébriques dont tous les conjugués complexes ont la valeur absolue $p^{i/2}$.

Sous les hypothèses et notations de (4.9), en rappelant que $(p,n)=1$, on a:

\begin{theorem}[5.1]
Si les conjectures de Weil sont vraies, alors les valeurs absolues de l'endomorphisme $F$ de ${}^k_nW_\ell$ sont des entiers algébriques, dont tous les conjugués complexes sont de valeur absolue $p^{k+1/2}$.
\end{theorem}

Supposons les conjectures de Weil.

\begin{lemma}[5.2, modulo Weil]
Soit $X$ un schéma lisse sur $\F_p$, réalisable comme ouvert dans un schéma projectif lisse $X^*$. Alors les valeurs absolues de l'endomorphisme $F^*$ de $\widetilde H^i(X,\Q_\ell)$ sont des entiers algébriques de valeur absolue $p^{i/2}$.
\end{lemma}

L'application naturelle de $H_c^i(X,\Q_\ell)$ vers $H^i(X,\Q_\ell)$ se factorise par $H^i(X^*,\Q_\ell)$:
\[
H_c^i(X,\Q_\ell)\to H^i(X^*,\Q_\ell)\to H^i(X,\Q_\ell),
\]
de sorte que, comme $\Gal(\overline\F_p/\F_p)$-module, $\widetilde H^i(X,\Q_\ell)$ est quotient d'un sous-objet de $H^i(X^*,\Q_\ell)$. \QED

\begin{lemma}[5.3, modulo Weil]
Soient $S$ un schéma lisse sur $\F_p$ et $f:A\to S$ un schéma abélien sur $S$. Supposons que $A$ puisse être réalisé comme ouvert dans un schéma projectif lisse $A^*$ sur $\F_p$. Alors le Frobenius géométrique $F^*$ de $\widetilde H^i(S,R^jf_*\Q_\ell)$ a pour valeurs propres des entiers algébriques de valeur absolue $p^{i+j/2}$.
\end{lemma}

Soit $m$ un entier $>1$, et considérons les suites spectrales de Leray
\[
E:
E_2^{ij}=H^i(\overline S,R^jf_*\Q_\ell)\Longrightarrow H^{i+j}(\overline A,\Q_\ell),
\qquad
{}_cE:
{}_cE_2^{ij}=H_c^i(\overline S,R^jf_*\Q_\ell)\Longrightarrow H_c^{i+j}(\overline A,\Q_\ell).
\]
L'endomorphisme de multiplication par $m$, $\psi_m=m\cdot1_A$, définit des endomorphismes de $E$ et de ${}_cE$ qui s'insèrent dans un diagramme commutatif
\[
\begin{array}{ccc}
{}_cE&\longrightarrow&E\\
\psi_m^*\downarrow&&\psi_m^*\downarrow\\
{}_cE&\longrightarrow&E.
\end{array}
\]
$\psi_m^*$ agit sur $R^jf_*\Q_\ell$ par multiplication par $m^j$, donc $\psi_m^*$ agit par multiplication par $m^j$ sur les termes ${}_cE_r^{ij}$ et $E_r^{ij}$ de ${}_cE$ et de $E$. Les applications $d_r$ pour $r\ge2$ commutent à $\psi_m^*$ et envoient $E_r^{ij}$, respectivement ${}_cE_r^{ij}$, dans $E_r^{i'j'}$, respectivement ${}_cE_r^{i'j'}$, avec $j\ne j'$. Elles sont nulles, et $E_2^{ij}$, respectivement ${}_cE_2^{ij}$, s'identifie à un sous-espace de $H^{i+j}(\overline A,\Q_\ell)$, respectivement de $H_c^{i+j}(\overline A,\Q_\ell)$, où $\psi_m^*=m^j$. Par conséquent, $\widetilde H^i(S,R^jf_*\Q_\ell)$ s'identifie, comme module galoisien, à un sous-espace de $\widetilde H^{i+j}(A,\Q_\ell)$ où $\psi_m^*=m^j$, et l'on applique (5.2). L'astuce utilisée ici est due à Lieberman. \QED

Soit $f_n:E\to M_n\otimes\F_p$ la courbe elliptique universelle sur $M_n\otimes\F_p$, et soit $f_{n,k}:E_k\to M_n\otimes\F_p$ son produit fibré $k$ fois avec lui-même. La formule de Künneth montre que le faisceau en $\Q_\ell$ $R^kf_{n,k*}\Q_\ell$ admet comme facteur direct la puissance tensorielle $k$-ième de $R^1f_{n*}\Q_\ell$; celle-ci contient à son tour comme facteur direct le faisceau en $\Q_\ell$ $\Sym^k(R^1f_{n*}\Q_\ell)$. Le théorème 5.1 résulte donc de (5.3) et du lemme suivant.

\begin{lemma}[5.4]
Le schéma $E_k$ est ouvert dans un schéma projectif lisse $E_k^*$ sur $\F_p$.
\end{lemma}

Soit $E^*$ le modèle minimal de Néron de $E$ sur $M_n^*\otimes\F_p$ (4.1). Le schéma $E^*$ est lisse et projectif sur $\F_p$. Comme $n\ge3$ et que les points d'ordre $n$ de $E$ forment un revêtement trivial de $M_n\otimes\F_p$, le modèle de Néron est semi-stable, cas a ou b$_m$ dans la classification de Néron. En particulier, la projection $f_n:E^*\to M_n^*$ n'a qu'un nombre fini de points non lisses, et, en ces points, $f_n$ est non dégénérée, présentant une singularité quadratique ordinaire.

Soit $E_k^{**}$ le produit fibré $k$-ième de $E^*$ sur $M_n^*$. Pour prouver (5.4), il suffit de résoudre les singularités de $E_k^{**}$ sans toucher l'ouvert $E_k$. On prouve d'abord:

\begin{lemma}[5.5]
Soit $V$ la sous-variété de l'espace affine sur un corps $k$, de coordonnées $(X_i)_{0\le i\le r}$, $(Y_i)_{0\le i\le r}$ et $(T_i)_{1\le i\le s}$, définie par l'équation
\[
X_0Y_0=X_1Y_1=\cdots=X_rY_r.
\]
Soit $\mathfrak m$ l'idéal de $\cO_V$ engendré par les monômes obtenus à partir du monôme $\prod_{i=0}^rX_i^i$ par une permutation des coordonnées qui respecte l'ensemble des paires $\{X_i,Y_i\}$, $0\le i\le r$. Alors $\mathfrak m=\cO_V$ hors du lieu singulier de $V$, et la variété $\widetilde V$ obtenue de $V$ par éclatement de $\mathfrak m$ est lisse sur $k$.
\end{lemma}

Le lieu singulier est le lieu où les quatre coordonnées $X_i,Y_i,X_j,Y_j$, pour $i\ne j$, s'annulent. L'ouvert affine de $\widetilde V$ défini par l'élément $\prod_{i=0}^rX_i^i$ de l'idéal $\mathfrak m$ de l'éclatement est le spectre de l'anneau régulier
\[
k[Y_0/X_1,X_0/X_1,X_1/X_2,\ldots,X_{r-1}/X_r,X_r,T_1,\ldots,T_s]
\]
(pour la preuve, remarquer que $X_i/X_{i+1}=Y_{i+1}/Y_i$), et 5.5 en résulte. \QED

On montre encore que, localement pour la topologie étale, les singularités de $E_k^{**}$ sont isomorphes à celles de $V$ pour $r=k-1$, et que cela permet de définir sur $E_k^{**}$ un idéal $\mathfrak m$ analogue à l'idéal $\mathfrak m$ de (5.5). En éclatant cet idéal, on obtient $E_k^*$. \QED

Une approximation du théorème suivant a été prouvée par Ihara [2].

\begin{theorem}[5.6]
Les conjectures de Weil impliquent la conjecture de Ramanujan.
\end{theorem}

On remarque d'abord que (5.1) est vrai pour $n=1$, car ${}^k_1W_\ell$ est le sous-module galoisien de ${}^k_mW_\ell$ invariant sous $\GL_2(\Z/(m))$. Sur ${}^k_1W_\ell$, $I_p^*$ induit l'identité, et (4.8) se réduit à
\[
1-T_pX+p^{k+1}X^2=(1-FX)(1-VX).
\]
Les endomorphismes $F$ et $V$ sont transposés l'un de l'autre, de sorte que
\[
\det(1-FX;{}^k_1W_\ell)=\det(1-VX;{}^k_1W_\ell).
\]
L'action de $T_p$ sur ${}^k_1W_\ell$ est induite par son action sur ${}^k_1W$, et est compatible avec la décomposition de ${}^k_1W\otimes\C$ en la somme de l'espace $S_{k+2}$ des formes paraboliques relatives à $\SL_2(\Z)$, de poids $k+2$, et de son conjugué complexe. De la propriété hermitienne de $T_p$, pour le produit intérieur de Peterson, et de (3.19), on déduit alors que
\[
\det(1-T_pX+p^{k+1}X^2;{}^k_1W_\ell)=\det(1-T_pX+p^{k+1}X^2;S_{k+2})^2,
\]
et
\[
\det(1-T_pX+p^{k+1}X^2;S_{k+2})^2=\det(1-FX;{}^k_1W_\ell)^2,
\]
d'où
\[
\det(1-T_pX+p^{k+1}X^2;S_{k+2})=\det(1-FX;{}^k_1W_\ell).
\tag{5.7}
\]
Reprenons les notations du paragraphe 1 et posons $k=10$. Par la théorie de Hecke et par (3.19), (5.7) se récrit
\[
H_p(X)=\det(1-FX;{}^{10}_1W_\ell),
\]
et l'on applique (5.1). \QED

On vérifie de la même manière que les conjectures de Weil impliquent la généralisation de Petersson de la conjecture de Ramanujan.

\section*{Références}
\begin{enumerate}[label={[\arabic*]},leftmargin=2.6em]
\item J. Igusa, \emph{Kroneckerian model of fields of elliptic modular functions}, Am. J. of Math. \textbf{81} (1959), 561--577.
\item Y. Ihara, \emph{Hecke polynomials as congruence zeta functions in elliptic modular case}, Ann. of Math., S.2 \textbf{85} (1967), 267--295.
\item J.-P. Jouanolou, exposés V et VI de SGA 5.
\item M. Kuga et G. Shimura, \emph{On the zeta function of a fibre variety whose fibers are abelian varieties}, Ann. of Math., S.2 \textbf{82} (1965), 478--539.
\item M. Raynaud, exposé XIII de SGA 1 et appendice (en préparation).
\item J.-P. Serre, \emph{Une interprétation des congruences relatives à la fonction $\tau$ de Ramanujan}, Sém. Delange--Pisot--Poitou, 1967/68, no. 14.
\item G. Shimura, \emph{Sur les intégrales attachées aux formes automorphes}, J. Math. Soc. of Japan \textbf{11} (1959), 291--311.
\item J. Tate, \emph{Courbes elliptiques: formulaire}, mis au goût du jour par P. Deligne, notes miméographiées par l'I.H.E.S.
\item J.-L. Verdier, \emph{Sur les intégrales attachées aux formes automorphes} (d'après G. Shimura), Sém. Bourbaki, février 1961, exp. 216.
\end{enumerate}

\paragraph{Sigles.}
EGA: \emph{Éléments de géométrie algébrique}, par A. Grothendieck et J. Dieudonné, Publ. Math. I.H.E.S.

SGA: \emph{Séminaire de géométrie algébrique du Bois-Marie}, notes miméographiées par l'I.H.E.S., en préparation chez North-Holland.

\end{document}
